% Russian resume, 1 page. Compiles with ANY engine -- no Overleaf config needed:
%  - Overleaf default (pdfLaTeX): uses T2A Cyrillic fonts automatically.
%  - XeLaTeX / LuaLaTeX / tectonic: uses fontspec + PT Serif (nicer look).
\documentclass[10pt,a4paper]{article}

\usepackage{iftex}
\ifpdftex
  \usepackage[utf8]{inputenc}
  \usepackage[T2A]{fontenc}
\else
  \usepackage{fontspec}
  % PT Serif ships with TeX Live / Overleaf and covers Cyrillic.
  \setmainfont{PT Serif}
\fi
\usepackage[english,russian]{babel}
\usepackage[a4paper,margin=1.4cm]{geometry}
\usepackage{enumitem}
\usepackage{titlesec}
\usepackage[colorlinks=true,urlcolor=blue]{hyperref}

\urlstyle{same}
\setlist[itemize]{leftmargin=1.2em,itemsep=0.36em,topsep=0.30em,parsep=0em,partopsep=0em}
\titleformat{\section}{\large\bfseries}{}{0em}{}[\vspace{-0.10em}\titlerule]
\titlespacing*{\section}{0pt}{0.74em}{0.30em}
\pagestyle{empty}
\setlength{\parindent}{0pt}
\setlength{\parskip}{0pt}
\emergencystretch=2em

\newcommand{\contactsep}{\enspace\textbar\enspace}

\begin{document}
\fontsize{9.8pt}{11.4pt}\selectfont

\begin{center}
{\LARGE \textbf{Артём Сиднев}}\\[0.10em]
{\normalsize \textbf{Software Engineer} \textbullet{} Java\,/\,Kotlin\,/\,Go \textbullet{} \textasciitilde4 года}\\[0.12em]
{\footnotesize Spring Boot \textbullet{} PostgreSQL \textbullet{} ClickHouse \textbullet{} Kafka \textbullet{} gRPC \textbullet{} Kubernetes}\\[0.20em]
{\footnotesize Москва \textbullet{} удалёнка\,/\,гибрид \textbullet{} открыт к релокации}\\[0.20em]
{\footnotesize
\href{tel:+79109692922}{+7 910 969 29 22}
\contactsep \href{mailto:a.sidnevart@gmail.com}{a.sidnevart@gmail.com}
\contactsep \href{https://t.me/sidnevart}{Telegram}
\contactsep \href{https://www.linkedin.com/in/artem-sidnev-b109ab401/}{LinkedIn}
\contactsep \href{https://github.com/sidnevart}{GitHub}
\contactsep \href{https://sidnevart.github.io/personal-web-page/}{Portfolio}
}
\end{center}

\vspace{-0.14em}

\section*{Профиль}
Backend \& Platform Engineer (Java, Kotlin, Go) с 4-летним опытом разработки высоконагруженных систем и Big Data пайплайнов. Технически возглавлял развитие Customer Data Platform (CDP), спроектировав архитектуру сбора данных, что увеличило выручку на \textasciitilde14 млн руб./мес.\ и ускорило процессы в 40 раз. Имею практический опыт внедрения AI-инструментов (LLM, RAG) в инженерные процессы. Специализируюсь на создании отказоустойчивых решений (ClickHouse, Kafka, gRPC) с фокусом на измеримый бизнес-результат.

\section*{Ключевые навыки}
\textbf{Языки:} Java, Kotlin, Python, Go, SQL \\
\textbf{Технологии:} Spring Boot, gRPC, REST, микросервисы, PostgreSQL, ClickHouse, CHaaS (ClickHouse-as-a-Service), Kafka, Redis, S3 \\
\textbf{Инструменты:} Docker, Kubernetes, GitLab CI/CD, Prometheus, Grafana, OpenTelemetry, nginx

\section*{Опыт}

\textbf{Software Engineer \textbar{} Tinkoff Bank} \hfill \textit{окт 2025 -- н.в.}\\[0.18em]
{\footnotesize\textit{Customer Data Platform (CDP): управление маркетинговыми аудиториями для кэшбэк-офферов (AdTech внутри банка)~--- отбор клиентов под партнёрские кампании по транзакционным, клиентским и партнёрским данным.}}
\begin{itemize}
  \item Ускорил сбор аудиторий \textbf{в 40--50 раз} (с 6 мин до 1--10 с), принеся \textbf{\textasciitilde14 млн руб./мес.} дополнительной выручки. Единолично спроектировал архитектуру DAG-пайплайнов на Kotlin, обеспечив мгновенную обработку данных любого объёма через CHaaS.
  \item \textbf{Технический лид интеграции CDP} (2 команды, 9 инженеров): инициировал RFC и спроектировал двухуровневую архитектуру~--- control-plane (\textbf{Kotlin\,/\,Spring Boot}) и data-plane (\textbf{Go}, gRPC, S3). Консолидировал смежную команду и сервис в контур CDP, ускорил доставку данных \textbf{в 32 раза} (с 8 ч до 15 мин) и достиг \textbf{99.9\%} надёжности процессов за счёт идемпотентности и retry\,/\,fallback.
  \item Архитектурно реализовал высоконагруженную обработку \textbf{\textasciitilde30 млн транзакций/день} через Kafka и ClickHouse с историей 1.5 года; система таргетинга на базе этого пайплайна принесла \textbf{\textasciitilde15 млн руб. годовой прибыли}.
  \item Сократил время выполнения критических сценариев \textbf{в 5 раз}, построив математическую модель задержек и сквозной мониторинг (observability) пайплайнов; оптимизировал критический путь миграцией вычислений в CHaaS.
  \item Экономлю менеджерам \textbf{до 60 рабочих часов ежемесячно} за счёт автоматизации проверок данных и настройки регионов; ускорил основные GitLab CI/CD пайплайны тестирования на 20\%.
\end{itemize}

\vspace{0.55em}
\textbf{Software Engineer \textbar{} Клиент (NDA)~--- рынок недвижимости} \hfill \textit{авг 2024 -- авг 2025}\\[0.18em]
{\footnotesize\textit{Аналитическая платформа для рынка жилой недвижимости.}}
\begin{itemize}
  \item Создал аналитическую платформу для рынка недвижимости: парсинг торгов и CIAN, геокодинг, каталог объектов с поиском, Telegram-бот и админ-панель, расчёт экономики объекта (Python, FastAPI, PostgreSQL, Redis).
  \item Сократил первичную оценку объекта \textbf{до \textasciitilde10 мин}, анализ района~--- до 7--10 мин: риэлторы и инвесторы видят цену, доходность, денежный поток и срок окупаемости без ручной сверки нескольких источников.
\end{itemize}

\vspace{0.55em}
\textbf{Freelance Software Engineer \textbar{} Клиентские проекты} \hfill \textit{май 2022 -- авг 2024}\\[0.18em]
{\footnotesize\textit{Разработка сервисов бронирования, лояльности и автоматизации для малого и сервисного бизнеса~--- от MVP до продакшена, единоличная ответственность за стек и поставку.}}
\begin{itemize}
  \item Разработал за \textbf{\textasciitilde1 месяц} в одиночку MVP, который помог заказчику привлечь инвестиции: Telegram Mini App для бронирования услуг~--- подбор свободных слотов и бронь в 2 действия вместо звонка, кабинеты клиента, партнёра и админа. Стек: Python (FastAPI), Go-сервис фоновой обработки и уведомлений, PostgreSQL, Redis, React; деплой Docker\,/\,nginx, CI/CD.
  \item Для бильярдного клуба перевёл запись из телефонных звонков в Telegram Mini App (\textbf{+25\% конверсии}): интеграция с YClients, система матчей с рейтингом \textbf{Elo} и идемпотентными бонусами лояльности~--- \textbf{экономия \textasciitilde15 ч/мес} администратора.
\end{itemize}

\section*{Образование}
\textbf{РосНОУ} \textbar{} Бакалавриат, Computer Science \textbar{} \textit{ожидаемое окончание: 2028}\\[0.05em]

\end{document}
